% 数学ノート集・共通テンプレート v2（積分論2026年度第1部の版面を参考）
% A5判。章内で定義・定理・例・注意等の番号を共通化。
% コンパイル： platex -interaction=nonstopmode math_note_template_v2.tex (2回)
%              dvipdfmx math_note_template_v2.dvi
\documentclass[a5paper,10pt,onecolumn]{jsarticle}
\usepackage[dvipdfmx,top=22mm,bottom=20mm,inner=19mm,outer=18mm,headheight=13pt,headsep=6mm]{geometry}
\usepackage{amsmath,amssymb,amsthm,mathtools}
\usepackage{enumitem,fancyhdr,titlesec}
\usepackage[dvipdfmx,bookmarks=true,colorlinks=false]{hyperref}
\usepackage{pxjahyper}

% ----- 差し替える文書情報 -----
\newcommand{\NoteTitle}{Feynman--Kacの公式とKolmogorov方程式}
\newcommand{\NoteSeries}{数学ノート集}

% ----- 文字・組版 -----
% 参考PDFの質素な明朝＋Computer Modernの印象を引き継ぎながら，
% A5判で窮屈にならないよう本文の行間を少し確保する．
\linespread{1.13}
\setlength{\parindent}{1zw}
\setlength{\parskip}{0pt}
\setlength{\abovedisplayskip}{6pt plus 2pt minus 2pt}
\setlength{\belowdisplayskip}{6pt plus 2pt minus 2pt}
\setlength{\abovedisplayshortskip}{3pt plus 2pt minus 1pt}
\setlength{\belowdisplayshortskip}{4pt plus 2pt minus 1pt}
\setlist[enumerate]{topsep=3pt,itemsep=1pt,parsep=0pt,leftmargin=2.4zw}
\titleformat{\section}{\normalfont\normalsize\bfseries}{\thesection}{0.8zw}{}
\titlespacing*{\section}{0pt}{2.0ex plus .4ex}{.9ex}
\titleformat{\subsection}{\normalfont\normalsize\bfseries}{\thesubsection}{0.7zw}{}
\titlespacing*{\subsection}{0pt}{1.5ex plus .3ex}{.6ex}
\raggedbottom

% ----- 柱とページ番号：ページ上部、罫線なし -----
\fancypagestyle{mathnote}{%
  \fancyhf{}%
  \fancyhead[L]{\footnotesize \NoteSeries\ （\NoteTitle）}%
  \fancyhead[R]{\footnotesize \thepage}%
  \renewcommand{\headrulewidth}{0pt}%
  \renewcommand{\footrulewidth}{0pt}%
}
\pagestyle{mathnote}

% ----- 節内通し番号の定義・定理群 -----
\newtheoremstyle{mathnotetheorem}
 {4.5pt}{4.5pt}{\normalfont}{}%
 {\bfseries}{}{0.65zw}%
 {\thmname{#1}\thmnumber{ #2}\thmnote{（#3）}}
\theoremstyle{mathnotetheorem}
\newtheorem{theorem}{定理}[section]
\newtheorem{definition}[theorem]{定義}
\newtheorem{lemma}[theorem]{補題}
\newtheorem{proposition}[theorem]{命題}
\newtheorem{corollary}[theorem]{系}
\newtheorem{example}[theorem]{例}
\newtheorem{remark}[theorem]{注意}
\renewcommand{\proofname}{証明}
\renewcommand{\qedsymbol}{$\square$}
\makeatletter
\renewenvironment{proof}[1][\proofname]{%
  \par\pushQED{\qed}\normalfont\topsep4\p@\@plus2\p@\relax
  \trivlist\item[\hskip\labelsep\bfseries #1]\ignorespaces
}{\popQED\endtrivlist\@endpefalse}
\makeatother

% ----- 使うノートで適宜拡張する記号 -----
\newcommand{\R}{\mathbb{R}}
\newcommand{\norm}[1]{\lVert #1\rVert}
\newcommand{\abs}[1]{\lvert #1\rvert}

\begin{document}
\section*{読了後にできること}
本稿を読むための前提は，Brown運動，熱方程式，Feynman--Kac公式．
本稿の具体的な到達課題は，期待値として与えられた関数を計算し，対応する偏微分方程式を検証する．

\section{拡散過程と生成作用素}

$B$ を $m$ 次元標準Brown運動とし，
$b:\mathbb R^d\to\mathbb R^d$，
$\sigma:\mathbb R^d\to\mathbb R^{d\times m}$
を有界で大域Lipschitz連続な写像とする．
このとき任意の初期値 $x\in\mathbb R^d$ に対して
\begin{equation}\label{eq:diffusion}
 X_t^x=x+\int_0^t b(X_s^x)\,ds
       +\int_0^t\sigma(X_s^x)\,dB_s
\end{equation}
は一意な連続適合解を持つ．

\begin{definition}
$a(x)=\sigma(x)\sigma(x)^{\mathsf T}$ と置く．
$C^2$ 級関数 $f:\mathbb R^d\to\mathbb R$ に対して
\[
 Lf(x)=\sum_{i=1}^d b_i(x)\partial_if(x)
 +\frac12\sum_{i,j=1}^d a_{ij}(x)\partial_{ij}f(x)
\]
を定める．これを拡散過程の\emph{生成作用素}という．
\end{definition}

\begin{proposition}[It\^oの公式による生成作用素]\label{prop:generator}
$f\in C_b^2(\mathbb R^d)$ とし，
その1階・2階導関数は有界かつ連続であるとする．
このとき
\[
 f(X_t^x)-f(x)-\int_0^tLf(X_s^x)\,ds
\]
は二乗可積分なマルチンゲールである．
\end{proposition}
\begin{proof}
多次元It\^oの公式により
\begin{align*}
 f(X_t^x)-f(x)
 &=\int_0^t Lf(X_s^x)\,ds\\
 &\quad+\sum_{i=1}^d\sum_{k=1}^m
   \int_0^t\partial_if(X_s^x)\sigma_{ik}(X_s^x)\,dB_s^k.
\end{align*}
係数と $\nabla f$ は有界なので，
最後の確率積分の被積分過程は二乗可積分である．
従って確率積分は真の二乗可積分マルチンゲールとなる．
\end{proof}

\begin{proposition}[生成作用素の極限表示]\label{prop:limit-generator}
$f\in C_b^2(\mathbb R^d)$ とし，その偏導関数は有界で連続，
$Lf$ は連続とする．
$P_tf(x)=E[f(X_t^x)]$ と置くと
\[
 \lim_{t\downarrow0}\frac{P_tf(x)-f(x)}{t}
 =Lf(x)
\]
が任意の $x\in\mathbb R^d$ で成立する．
\end{proposition}
\begin{proof}
命題\ref{prop:generator}で期待値を取ると
\[
 P_tf(x)-f(x)=E\int_0^tLf(X_s^x)\,ds.
\]
係数が有界なので
\[
 E|X_s^x-x|^2
 \le2sE\int_0^s|b(X_u^x)|^2du
   +2E\int_0^s|\sigma(X_u^x)|_{\mathrm{HS}}^2du
 \le Cs
\]
となる．従って $X_s^x\to x$ は確率収束である．
$Lf$ は連続で有界なので
$E[Lf(X_s^x)]\to Lf(x)$ が従う．
時間平均を取って
\[
 \frac1t\int_0^tE[Lf(X_s^x)]\,ds\to Lf(x)
\]
を得る．
\end{proof}

\section{Markov半群}

\begin{definition}
有界Borel関数 $f$ に対し
\[
 (P_tf)(x)=E[f(X_t^x)]
\]
と定める．$(P_t)_{t\ge0}$ を拡散過程の
\emph{Markov半群}という．
\end{definition}

\begin{theorem}[Markov性と半群性]\label{thm:semigroup}
任意の有界Borel関数 $f$ と $s,t\ge0$ について
\[
 P_{s+t}f=P_s(P_tf)
\]
が成立する．また $P_t1=1$，
$f\ge0$ なら $P_tf\ge0$，
$\norm{P_tf}_\infty\le\norm f_\infty$ である．
\end{theorem}
\begin{proof}
時刻 $s$ までのBrown運動の情報を $\mathcal F_s$ とする．
$s$ 以降のBrown運動の増分
$\widetilde B_u=B_{s+u}-B_s$ は
$\mathcal F_s$ と独立な標準Brown運動である．
経路ごとの一意性により，
時刻 $s$ 以降の解は $X_s^x$ を初期値とし
$\widetilde B$ で駆動される同じ確率微分方程式の解となる．
従って
\[
 E[f(X_{s+t}^x)\mid\mathcal F_s]
 =(P_tf)(X_s^x).
\]
期待値を取れば $P_{s+t}f=P_s(P_tf)$ を得る．
残る主張は $P_t$ が確率測度に関する期待値として
定義されていることから従う．
\end{proof}

\begin{proposition}[初期値の連続性]\label{prop:feller}
任意の有界連続関数 $f$ に対して
$P_tf$ は有界連続関数である．
\end{proposition}
\begin{proof}
大域Lipschitz条件のもとでは，
初期値 $x,y$ に対応する解について
\[
 E\sup_{s\le t}|X_s^x-X_s^y|^2
 \le C_t|x-y|^2
\]
がGronwallの不等式から従う．
$x_n\to x$ なら $X_t^{x_n}\to X_t^x$ は確率収束する．
$f$ は連続かつ有界なので
$f(X_t^{x_n})\to f(X_t^x)$ は確率収束し，
一様有界性によって期待値も収束する．
従って $P_tf(x_n)\to P_tf(x)$ を得る．
\end{proof}

\section{後退Kolmogorov方程式}

\begin{theorem}[後退方程式の古典解の表示]
\label{thm:backward}
$T>0$ とし，
$u\in C^{1,2}([0,T]\times\mathbb R^d)$ が有界で，
$u_t,\nabla u,D^2u$ も有界であるとする．
\[
 \partial_tu(t,x)+Lu(t,x)=0,\quad t<T,
 \qquad u(T,x)=g(x)
\]
を満たすなら
\[
 u(t,x)=E[g(X_{T-t}^x)]
 \qquad(0\le t\le T)
\]
が成立する．従ってこのクラスの古典解は一意である．
\end{theorem}
\begin{proof}
固定した $(t,x)$ について，
$Y_s=X_s^x$（$0\le s\le T-t$）と置く．
It\^oの公式を $s\mapsto u(t+s,Y_s)$ に適用すると
\begin{align*}
 u(t+s,Y_s)-u(t,x)
 &=\int_0^s(\partial_tu+Lu)(t+r,Y_r)\,dr\\
 &\quad+\int_0^s\nabla u(t+r,Y_r)
     \sigma(Y_r)\,dB_r.
\end{align*}
偏微分方程式により第1項は零である．
第2項の被積分過程は有界なので期待値は零となる．
$s=T-t$ と置いて期待値を取れば
$u(t,x)=E[u(T,Y_{T-t})]=E[g(X_{T-t}^x)]$
を得る．右辺は $u$ の選択に依存しないから一意性も従う．
\end{proof}

\section{前進Kolmogorov方程式}

\begin{proposition}[法則の弱方程式]\label{prop:forward-weak}
$\mu_t$ を $X_t^x$ の法則とする．
任意の $f\in C_c^\infty(\mathbb R^d)$ に対して
\[
 \int f(y)\,\mu_t(dy)-f(x)
 =\int_0^t\int Lf(y)\,\mu_s(dy)\,ds.
\]
\end{proposition}
\begin{proof}
命題\ref{prop:generator}のマルチンゲールの期待値を取ると
\[
 E[f(X_t^x)]-f(x)
 =\int_0^tE[Lf(X_s^x)]\,ds.
\]
各期待値を法則 $\mu_t$ による積分に書き換えればよい．
\end{proof}

\begin{corollary}[Fokker--Planck方程式]\label{cor:fokker}
$\mu_t$ がLebesgue密度 $p(t,y)$ を持ち，
係数と密度が必要な微分・部分積分を正当化する程度に
滑らかであると仮定する．このとき分布の意味で
\[
 \partial_tp(t,y)
 =-\sum_{i=1}^d\partial_{y_i}(b_i(y)p(t,y))
  +\frac12\sum_{i,j=1}^d
    \partial_{y_i}\partial_{y_j}(a_{ij}(y)p(t,y))
\]
が成立する．
\end{corollary}
\begin{proof}
命題\ref{prop:forward-weak}の両辺を $t$ で微分すると
\[
 \frac d{dt}\int f(y)p(t,y)\,dy
 =\int\left(
 \sum_i b_i\partial_if+\frac12\sum_{i,j}a_{ij}\partial_{ij}f
 \right)p(t,y)\,dy.
\]
$f$ はコンパクト台を持つので，
各項で部分積分を1回または2回行えば
右辺は表示された微分作用素を $p$ に適用したものとの
双対対合となる．
任意の $f\in C_c^\infty$ について成立するため，
分布の意味で方程式を得る．
\end{proof}

\begin{remark}
この系は確率過程の法則が既に密度を持つ場合の主張である．
一般の係数から密度の存在や滑らかさが自動的に従うとは
主張していない．
\end{remark}

\section{Dynkinの公式と停止時刻}

\begin{theorem}[Dynkinの公式]\label{thm:dynkin}
$\tau\le T$ を停止時刻，
$f\in C_b^2(\mathbb R^d)$ とし，
導関数も有界とする．このとき
\[
 E[f(X_\tau^x)]-f(x)
 =E\int_0^\tau Lf(X_s^x)\,ds.
\]
\end{theorem}
\begin{proof}
命題\ref{prop:generator}のマルチンゲールを
停止時刻 $\tau$ で止める．
積分過程は二乗可積分で $\tau$ は有界だから，
停止定理により停止したマルチンゲールの期待値は零である．
従って
\[
 E\left[
 f(X_\tau^x)-f(x)-\int_0^\tau Lf(X_s^x)\,ds
 \right]=0
\]
となり結論を得る．
\end{proof}

\begin{example}[区間からの退出時刻]\label{ex:exit}
1次元Brown運動を $X_t=x+B_t$，
$-a<x<a$，$a>0$ とし，
$\tau=\inf\{t\ge0:|X_t|=a\}$ と置く．
このとき $\tau<\infty$ が概至る所成立し，
\[
 E[\tau]=a^2-x^2,\qquad
 P(X_\tau=a)=\frac{x+a}{2a}.
\]
\end{example}
\begin{proof}
$f(y)=a^2-y^2$ と置く．
$|y|\le a$ の範囲では $f(y)\ge0$ であり，
Brown運動の生成作用素は $L=\frac12d^2/dy^2$
だから $Lf=-1$ となる．
$\tau_n=\tau\wedge n$ とし，
停止区間上だけで必要な偏導関数が有界な
$C_b^2$ 関数へ $f$ を切断して
定理\ref{thm:dynkin}を適用すると
\[
 E[a^2-X_{\tau_n}^2]
 =a^2-x^2-E[\tau_n].
 \tag{$*$}
\]
左辺は非負なので $E\tau_n\le a^2-x^2$ である．
単調収束定理により
$E\tau\le a^2-x^2<\infty$ となり，
特に $\tau<\infty$ が概至る所成立する．
Brown運動の連続性から $X_\tau\in\{-a,a\}$ である．
$0\le a^2-X_{\tau_n}^2\le a^2$ を用いて
$(*)$ に有界収束定理を適用すれば
$E\tau=a^2-x^2$ である．

次に $g(y)=(y+a)/(2a)$ と置くと $Lg=0$ である．
同じ停止論法から
$Eg(X_{\tau_n})=g(x)$ となる．
有界収束定理を使って $n\to\infty$ とすると
$Eg(X_\tau)=g(x)$ を得る．
$g(-a)=0$，$g(a)=1$ だから
$P(X_\tau=a)=(x+a)/(2a)$ となる．
\end{proof}

\section{Feynman--Kacの公式}

\begin{theorem}[Feynman--Kacの表示]\label{thm:fk}
$T>0$ とし，$b,\sigma$ は第1節の条件を満たすとする．
$V,F:[0,T]\times\mathbb R^d\to\mathbb R$ を
有界連続関数，$g:\mathbb R^d\to\mathbb R$ を
有界連続関数とする．
$u\in C^{1,2}([0,T]\times\mathbb R^d)$ が有界で，
$u_t,\nabla u,D^2u$ も有界であり，
\[
 \partial_tu+Lu-Vu+F=0,\qquad
 u(T,x)=g(x)
\]
を満たすなら，各 $(t,x)$ に対して
\begin{align}
 u(t,x)
 &=E\left[
 e^{-\int_t^TV(s,X_s^{t,x})\,ds}
    g(X_T^{t,x})\right]\notag\\
 &\quad+
 E\left[\int_t^T
  e^{-\int_t^sV(r,X_r^{t,x})\,dr}
  F(s,X_s^{t,x})\,ds\right].
 \label{eq:fk}
\end{align}
ここで $X_s^{t,x}$ は時刻 $t$ に $x$ から出発する
\eqref{eq:diffusion} の解とする．
従って上記の正則性クラスにおける古典解は一意である．
\end{theorem}
\begin{proof}
固定した $(t,x)$ を取り，
$s\in[t,T]$ に対して
\[
 A_s=\exp\!\left(-\int_t^sV(r,X_r^{t,x})\,dr\right)
\]
と置く．$A_t=1$ であり，
$A$ は有限変動過程で
$dA_s=-V(s,X_s^{t,x})A_s\,ds$ を満たす．
It\^oの積の公式を
$A_su(s,X_s^{t,x})$ に適用する．
$A$ は有限変動なので二次共変動の項は零であり，
\begin{align*}
 d(A_su(s,X_s))
 &=A_s\left(
 (\partial_su+Lu-Vu)(s,X_s)\,ds
 +\nabla u(s,X_s)\sigma(X_s)\,dB_s
 \right)\\
 &=-A_sF(s,X_s)\,ds+
 A_s\nabla u(s,X_s)\sigma(X_s)\,dB_s.
\end{align*}
$V$ は有界なので $A_s$ は決定論的な定数で有界であり，
$\nabla u,\sigma$ も有界である．
従って最後のIt\^o積分は二乗可積分マルチンゲールで，
期待値は零である．
$s=t$ から $s=T$ まで積分し，期待値を取ると
\[
 E[A_Tu(T,X_T)]-u(t,x)
 =-E\int_t^TA_sF(s,X_s)\,ds.
\]
$u(T,\cdot)=g$ を代入して移項すれば
\eqref{eq:fk} を得る．
右辺は $u$ の選択によらないため一意性も従う．
\end{proof}

\begin{remark}
定理\ref{thm:fk}は，既に存在する古典解の確率的表示と
一意性を証明する定理である．
一般の $b,\sigma,V,F,g$ のもとで右辺が
$C^{1,2}$ の古典解になることまでは，
追加の楕円型・放物型正則性理論なしに主張しない．
\end{remark}

\section{熱方程式とGaussian核}

\begin{theorem}[熱方程式の明示解]\label{thm:heat}
$g\in C^2(\mathbb R)$ とし，
$g,g',g''$ が有界で連続であるとする．
\[
 u(t,x)=\int_{\mathbb R}
 g(x+y)\frac{1}{\sqrt{2\pi(T-t)}}
 \exp\!\left(-\frac{y^2}{2(T-t)}\right)\,dy
 \quad(0\le t<T)
\]
および $u(T,x)=g(x)$ と置く．
このとき $u\in C^{1,2}([0,T)\times\mathbb R)$ で
\[
 \partial_tu+\frac12\partial_{xx}u=0
\]
を満たす．
また $t\uparrow T$ で $u(t,x)\to g(x)$ が成り立つ．
\end{theorem}
\begin{proof}
$h=T-t>0$ と置き，
$p_h(y)=(2\pi h)^{-1/2}e^{-y^2/(2h)}$ とする．
$g,g',g''$ は有界なのでLebesgueの収束定理によって
\[
 u_x(t,x)=\int g'(x+y)p_h(y)\,dy,\quad
 u_{xx}(t,x)=\int g''(x+y)p_h(y)\,dy
\]
と微分できる．
Gaussian核について
$\partial_hp_h=\frac12\partial_{yy}p_h$ が直接計算で従う．
積分と $h$ による微分を交換して
\[
 \partial_hu
 =\frac12\int g(x+y)\partial_{yy}p_h(y)\,dy.
\]
$g,g'$ が有界でGaussian核が無限遠で急減少するため，
部分積分を二回行えば
$\partial_hu=\frac12\int g''(x+y)p_h(y)\,dy
=\frac12u_{xx}$ である．
$\partial_t=-\partial_h$ なので
$u_t+\frac12u_{xx}=0$ を得る．
最後に $u(t,x)=E[g(x+\sqrt h\,Z)]$，
$Z\sim N(0,1)$ と書き，
$h\downarrow0$ で有界収束定理を適用すれば
$u(t,x)\to g(x)$ となる．
\end{proof}

\begin{corollary}
定理\ref{thm:heat}の解は
\[
 u(t,x)=E[g(x+B_{T-t})]
\]
と表せる．
\end{corollary}
\begin{proof}
$B_{T-t}\sim N(0,T-t)$ なので，
その密度を用いた期待値の定義は
定理\ref{thm:heat}の積分表示と一致する．
\end{proof}

\begin{remark}
確率微分方程式から生成作用素が定まり，
その随伴作用素が法則の時間発展を，
生成作用素自身が期待値の時間発展を支配する．
Feynman--Kacの公式は，この対応を
ポテンシャル項・外力項を含む方程式へ拡張する．
\end{remark}

\section{到達課題の完全解答}
\noindent\textbf{熱方程式の確率表示．}
\begin{proposition}
標準Brown運動 $B_t$ に対し
\[
u(t,x)=E[(x+B_t)^2]
\]
と置くと $u(t,x)=x^2+t$ であり，
\[
\partial_tu=\frac12\partial_{xx}u,\qquad
u(0,x)=x^2
\]
を満たす．
\end{proposition}
\begin{proof}
$E[B_t]=0$，$E[B_t^2]=t$ だから
\[
E[(x+B_t)^2]=x^2+2xE[B_t]+E[B_t^2]=x^2+t.
\]
従って $\partial_tu=1$，
$\partial_{xx}u=2$ であり方程式が成立する．
$t=0$ では $B_0=0$ だから初期条件も成立する．
この計算はFeynman--Kac公式の単純な例である．
\end{proof}
\noindent\textbf{到達確認．}
$g(x)=x$ の場合に $E[g(x+B_t)]$ を求めよ．
\par\noindent\textbf{解答．}
$E[x+B_t]=x$ であり，時間に依らない解となる．
\end{document}
